\documentclass[12pt, letterpaper]{article}

\setlength{\topmargin}{-1.75cm} \setlength{\textheight}{22.5cm}
\setlength{\oddsidemargin}{0.25cm}
\setlength{\evensidemargin}{0.25cm} \setlength{\textwidth}{16.2cm}

% symbols 
\usepackage{fdsymbol}
\usepackage{stmaryrd}
\usepackage{pifont} % check/cross marks 
\newcommand{\cmark}{\ding{51}}%
\newcommand{\xmark}{\ding{55}}%
%\usepackage{amssymb}
\usepackage[ruled,linesnumbered]{algorithm2e}
\usepackage{graphicx}
\usepackage[linguistics]{forest}
% for automata 
\usepackage{tikz} 
\usetikzlibrary{automata,positioning}


\usepackage{palatino, url, multicol} % for multiple columns

\usepackage{pictex}
%% in the .pictex output of xfig, there is command \colo
%% however the old version of pictex may not define this
%% so we define color here as empty
\def \color#1]#2{}

\usepackage{fancyhdr}
\pagestyle{fancy}

\newcommand{\hwnumber} {
    5
}

\newcommand{\hwtitle} {
  Push Down Automata.  Turing Machine.
}

\newcommand{\duedate}{
    11:59 pm, Wed November 19.
}

\date{}

\begin{document}

\newcommand{\hide}[1]{}
\newcommand{\otherquestions}[1]{}
\newcommand{\set}[1]{\{#1\}}
\newcommand{\pg}[1]{{\tt #1}}
\newtheorem{definition}{Definition}
\newcommand{\emptyclause}{\Box}
\newcommand{\keyword}[1]{{\bf #1}}

\newcommand{\ee}[1] {
  \begin{enumerate}
    #1
  \end{enumerate}
}
%% \noindent test \hfill test

\newcommand{\ie}[1] {
  \begin{enumerate}
    #1
  \end{enumerate}
}

\newcommand{\kleene}{\wedge}

\title{{\bf Task \hwnumber.} \hwtitle}

\maketitle

\input{hwHeader}

 \noindent {\bf Submit PDF file (and Latex if you use it) of your solution to Canvas by
  \duedate.}

\ee{

	% basics 
    \item (15 - basics) (Push down automata).  Given the PDA (push down automaton) for the palindrome (see lecture notes L5.2). Write the precise tuple form for this automaton. Remember to define each element of the tuple.  

	% 80% 
	\item (10 - medium) (PDA) Design PDA for the language $L = \{0^n 1^n | n  \ge  1\}$ %  (change 0, 1 to a, b)

	% 80% 
	%\item (10 - medium)  (PDA) Design PDA for $L = \{ w c w^R   | w \in  \{a, b\}^* \}$. Recall that $w^R$ is the reverse of $w$. 

	%	90% 
	\item (5 - advanced) (Comprehensive). Write  a precise and general description of how a push down automaton can be constructed for a context free grammar $G$ such that the automaton accepts  same language represented/generated by $G$. 

		%Working backwards, write a few steps for proving the following theorem.  Ideally, you can reach the step how the PDA is constructed.
	
		%	{\bf Theorem}. For every context free language $L$, there exists a push down automaton $M$, $L$ is accepted by a PDA. 
		
		% A more ``normal'' statement is like "Every context free language is accepted by a PDA." You are expected to reformulate such statement into the statement above where all logical symbols are explicit and which helps produce a proof (and thus the statement is indeed true). 

	\item (12) (pumping lemma idea) Consider the CFG 
	
	$ E \rightarrow E - E$ \\
	$ E \rightarrow number$ \\ 

	and the parse tree for $6-7-8$. 

	\includegraphics[width=0.25\textwidth]{hw5-parseTree.png}

	We use the tree to show the intuitive ideas of pumping lemma. 
	\ee{
	  \item Produce a new parse tree by replacing the subtree rooted at E with green circle by the subtree rooted at the E with red circle. 
	  \item What is the yield of the new tree? Hint. Yield refers to the string/expression formed by the leaves of the tree in the order from the left to right. 
	  \item Why the new tree is still a valid parse tree with respect to the CFG?   
	}

    \item (8) (pumping lemma) Consider the language $L = \{ a^i b^i c^i ~|~ i \ge 0 \}$. Assume there exists CFG $G=(N, T, P, S)$ for $L$. Recall that $\phi(G)$ refers to the maximum number of terminals and non-terminals in the right side of any production rules.
    % defining a non-terminal symbol. 

   Let string $w = a^n b^n c^n$ where $n > \phi(G)^{|N|}$.  By pumping lemma, there exists $u, v, x, y, z$ such that $w = uvxyz$, either $v$ or $y$ is not an empty string, $|vxy| \le n$, and for any $k$, $u v^k x y^k z \in L$. 

   \ee{ 
		\item Write a proof for the statement: if $vxy$ does not contain $c$, then $uxz \in L$. 
		\item write a proof for $uxz \not \in L$. 
   } 
   % Consider the case that  string $vxy$ does not contain $c$. Show there is a contradiction. 
  \item (28) (Turing Machine) Consider the Turing machine $M$ in the picture below
     \begin{center}
		\includegraphics[width=0.6\textwidth]{hw5-TuringMachine.png}
	 \end{center}
	 \ee{
		\item What are the states of the machine? 
		\item What is the alphabet of the language that is accepted by this machine? 
		\item What are the halting states? (Halting states the nodes with double circle in the diagram.)
		\item  What is the initial state? 
		\item  Write the {\em initial configuration} of $M$ on string $abc$.  (see definition of initial configuration in the slides).
		\item  Write the computation from $(q_1, \triangleright \textvisiblespace \  \underline{d}ab)$ to $(q_3, \triangleright \textvisiblespace \ 
		dd\underline{d})$.
		\item what is the value (as a tuple) of $\delta(q_2, b)$? 
		% $\delta(q_2, b) = (q_3, d)$. 
		% Write the part of transition function $\delta$ of $M$ where $q_2$ is an input state by completing the following set $\{((q_2, b), (q_3,d)), ... \}$. Note the notion  $((q_2, b), (q_3,d))$ represents $\delta(q_2, b) = (q_3, d)$. 
	}
    \item (5) (Church-Turing Thesis) What is Church-Turing thesis? 
 \hide{ 
     \item (5) (Turing Machine) What are the differences between the following two definitions of $K$. 
    
     \ie{
    		\item Definition 1.  {\bf $K$} is a set of states. 
    		\item Definition 2. {\bf $K$} is a set. Each element of $K$ is called a {\bf state}. 
      }
 }

  \item (10) (Turing Machine) Give a precise definition for ``A Turing machine rejects a string'' by assuming that the machine has only two halting states $y$ and $n$.  To make it easy for  your writing to refer to the machine and string, you may give a name to each of them such  as $M$ and $w$ respectively. Thus the concept becomes ``A Turing machine $M$ rejects a string $w$.'' You may find the definition of a similar concept of ``accepts" in lecture note L6.1. 
%	You can directly use the concepts such as To refer to the details of a Turing machine, you may need to introduce $(K, \Sigma \cup \{\}, s, H, \sigma)$, and for easy reference of a string, you may use $w$. Thus the about concept would be "A Turing machine    $(K, \Sigma \cup \{\}, s, H, \sigma)$ rejects a string $w$." 
  \item (Halting problem). What is halting problem? Give a proof that halting problem is not decidable. You are expected to apply the precise definition of halting problem, decideable and etc. 
}
\end{document}

